\section{Referencial Teórico e Estado da Arte}
\label{sec:referencial}

A análise automatizada de riscos em contratações públicas articula fundamentos da Ciência da Computação (algoritmos em grafos e estruturas de dados), da Ciência das Redes Complexas e da regulação jurídica e antitruste das licitações. Esta seção consolida o arcabouço normativo, a literatura científica e as formulações matemáticas que sustentam os três eixos do estudo.

\subsection{Arcabouço Jurídico e Tipologias de Risco Licitatório}

As licitações públicas visam assegurar a isonomia entre competidores e selecionar a proposta mais vantajosa para a Administração, conforme os princípios expressos nos arts.~5º, 9º, inciso~I, alínea ``a'', e 11, incisos~I a III, da Lei nº 14.133/2021~\cite{brasil2021licitacoes}. O mesmo diploma impõe pesquisa de preços compatível com o mercado e consideração da economia de escala (art.~23, §§~1º e 4º), tipifica infrações administrativas graves (art.~155, incisos~VIII a XII), estabelece a dosimetria das sanções conforme o dano apurado (art.~156, §~1º, incisos~I e IV), autoriza a desconsideração da personalidade jurídica para alcançar sócios e empresas coligadas (art.~160) e inclui o art.~337-F no Código Penal para coibir a frustração do caráter competitivo.

Sob a ótica concorrencial, a Lei nº 12.529/2011 (art.~36, \textit{caput}, §~3º, incisos~I, alíneas ``c'' e ``d'', e II)~\cite{brasil2011concorrencia} enquadra como infrações à ordem econômica a divisão concertada de mercados e a combinação de preços ou abstenções em certames. Complementarmente, a Lei Anticorrupção (Lei nº 12.846/2013, art.~5º, inciso~IV)~\cite{brasil2013anticorrupcao} responsabiliza objetivamente pessoas jurídicas por atos lesivos em licitações, enquanto a Lei Complementar nº 123/2006 (art.~3º, §~4º, incisos~III a V)~\cite{brasil2006lc123} veda a fruição de tratamento diferenciado de microempresas (ME) e empresas de pequeno porte (EPP) quando há participação societária cruzada que exceda o teto legal. Já a Lei de Improbidade Administrativa (Lei nº 8.429/1992, com redação dada pela Lei nº 14.230/2021, art.~10 e art.~12, inciso~II)~\cite{brasil1992improbidade} exige demonstração de dolo específico e de perda patrimonial efetiva para responsabilização, o que fundamenta o cruzamento entre indícios estruturais e materialidade financeira. Por fim, a jurisprudência do TCU~\cite{tcu2011acordao2341,tcu2024acordao1798} consolida o entendimento de que a existência isolada de sócio em comum ou parentesco não configura ilícito por si só, exigindo a convergência de múltiplos indícios objetivos para justificar apuração aprofundada.

Na literatura econômica e nos guias da Organização para a Cooperação e Desenvolvimento Econômico (OCDE)~\cite{oecd2025guidelines} e do Conselho Administrativo de Defesa Econômica (CADE)~\cite{cade2019guia}, os arranjos restritivos à competição em compras públicas manifestam-se em quatro mecanismos principais: (i) \textit{cover bidding} (propostas de cobertura), em que licitantes submetem lances elevados para simular disputa e viabilizar a vitória de um parceiro~\cite{carbone2024bidrigging,aaltio2025complementary}; (ii) supressão de propostas, caracterizada pela abstenção coordenada de concorrentes aptos; (iii) \textit{bid rotation} (rodízio de vencedores), marcado pela alternância cíclica de vitórias ao longo de certames sucessivos~\cite{imhof2018screening,imhof2021detecting}; e (iv) alocação setorial ou territorial, na qual fornecedores dividem secretarias ou nichos específicos, elevando artificialmente a concentração orçamentária~\cite{wachs2019cartel,wachs2020corruption}. Como os participantes costumam cumprir formalmente os requisitos habilitatórios, a detecção preventiva depende de métricas objetivas e relacionais capazes de rastrear padrões anômalos em grandes volumes de registros~\cite{fazekas2016objective}.

\subsection{Modelagem em Redes Complexas para Auditoria Pública}

A Ciência das Redes permite traduzir bases administrativas relacionais em estruturas formais analisáveis algoritmicamente~\cite{newman2018networks}. Em contratações públicas municipais brasileiras, estudos recentes demonstram a eficácia de projeções empresa--empresa de co-disputa~\cite{lyra2021characterization} e de grafos heterogêneos que acoplam co-participação em certames a vínculos societários do QSA~\cite{villamil2024collusion,pompeu2025collusion}. Para capturar integralmente o ciclo das despesas municipais sem distorcer a semântica de cada etapa, este trabalho adota três representações complementares:

\begin{enumerate}
  \item \textbf{Grafo Heterogêneo Multicamada ($G_{\text{soc}}$):} definido como $G_{\text{soc}} = (V_{\text{soc}}, E_{\text{soc}}, \tau, \varphi, w)$, em que $\tau$ particiona $V_{\text{soc}}$ nas classes \textit{Órgão}, \textit{Empresa} e \textit{Sócio Real}, e $\varphi$ distingue arestas de participação, co-concorrência, propriedade societária (\textit{É\_Dono\_De}) e vínculo cadastral ou societário suspeito (\textit{Vinculo\_Suspeito}), permitindo arestas paralelas ponderadas por $w: E_{\text{soc}} \to \mathbb{R}_{>0}$.
  \item \textbf{Dígrafo Ponderado de Disputas ($G_{\text{dir}}$):} definido como $G_{\text{dir}} = (V_{\text{emp}}, E_{P \to V}, w_{\text{dir}})$, em que uma aresta direcionada $u \to v$ registra que a empresa perdedora $u$ foi superada pela vencedora $v$ em itens multi-participante. Em bases públicas como o SICOM~\cite{tcemg2025in01}, nas quais os lances perdedores frequentemente aparecem zerados nos registros de julgamento, a modelagem direcionada preserva a assimetria vitória/derrota e viabiliza a análise de reciprocidade temporal e supressão de desconto frente à cotação prévia~\cite{imhof2021detecting,wachs2019cartel}.
  \item \textbf{Grafo Bipartido Unidade--Credor ($G_{\text{bip}}$):} modelado como rede de dois modos $G_{\text{bip}} = (U \cup C, E_{\text{emp}}, w_{\text{fin}})$~\cite{borgatti1997network}, ligando unidades orçamentárias $u \in U$ aos credores $c \in C$ com peso igual ao valor líquido empenhado. Aplicações similares na América Latina evidenciam que a análise bipartida revela dependência crônica e fornecedores-ponte em administrações descentralizadas~\cite{herrera2020characterization,alcazar2026supplier}.
\end{enumerate}

\subsection{Algoritmos de Grafos e Métricas Aplicadas}

Para processar as três representações dentro de limites determinísticos de complexidade assintótica, emprega-se o seguinte conjunto de algoritmos e métricas:

\subsubsection{Redução por $1$-Core, Union-Find e Listagem de Triângulos $K_3$}
Dado o grafo heterogêneo $G_{\text{soc}}$, filtra-se o subgrafo induzido pelas arestas \textit{Vinculo\_Suspeito} acima do limiar de risco e suas respectivas arestas de propriedade. Em seguida, aplica-se a decomposição $k$-core com $k=1$~\cite{seidman1983network} via fila em busca em largura ($\mathcal{O}(|V| + |E|)$), podando vértices isolados. Sobre o subgrafo resultante, a estrutura \textit{Disjoint-Set Union} (Union-Find) com compressão de caminhos e união por posto~\cite{tarjan1975efficiency} extrai as componentes conexas em tempo quase linear $\mathcal{O}(|E|\,\alpha(|V|))$, onde $\alpha$ é a inversa da função de Ackermann. Por fim, para cada par suspeito $(a,b)$, interseccionam-se as listas de adjacência de sócios em tempo $\mathcal{O}(|E|\,d_{\max})$ para enumerar triângulos societários fechados ($K_3$: Sócio--Empresa$_A$--Empresa$_B$).

\subsubsection{Componentes Fortemente Conexas (Tarjan), Ciclos Curtos e Comunidades de Louvain}
No dígrafo $G_{\text{dir}}$, o algoritmo de Tarjan~\cite{tarjan1972depth} identifica todas as Componentes Fortemente Conexas (SCCs) em uma única travessia em profundidade com complexidade $\mathcal{O}(|V| + |E|)$. Enquanto relacionamentos acíclicos geram SCCs unitárias ($|C| = 1$), a alternância de vitórias entre concorrentes produz SCCs não triviais ($|C| \ge 2$). Para isolar alternâncias restritas a múltiplas licitações, executa-se uma busca em profundidade limitada para ciclos simples $C_3$ e $C_4$ ($\mathcal{O}((|V|+|E|)c)$, sendo $c$ o fator de ramificação local). Complementarmente, o algoritmo de Louvain~\cite{blondel2008fast} maximiza a modularidade $Q$ na projeção não direcionada de $G_{\text{dir}}$:
\begin{equation}
  Q = \frac{1}{2m} \sum_{x, y \in V_{\text{emp}}} \left[ A_{xy} - \frac{k_x k_y}{2m} \right] \delta(g_x, g_y),
\end{equation}
onde $A_{xy}$ é o peso da aresta, $k_x = \sum_y A_{xy}$, $m = \frac{1}{2}\sum_{xy} A_{xy}$, $g_x$ é a comunidade do vértice $x$ e $\delta$ é a função delta de Kronecker. Ressalta-se que valores elevados de $Q$ e a presença de ciclos $C_3/C_4$ indicam apenas segmentação de mercado ou disputa recorrente entre especialistas de um mesmo nicho, exigindo confronto com razões de preço e datas distintas de certames~\cite{imhof2018screening,aaltio2025complementary}. Para caracterização estrutural comparativa da rede de co-disputas no painel de síntese, utiliza-se também a centralidade de intermediação não ponderada $C_B(v)$~\cite{freeman1977set}, computada pelo algoritmo de Brandes~\cite{brandes2001faster} em tempo $\mathcal{O}(|V|\cdot|E|)$.

\subsubsection{Concentração Orçamentária (HHI, Entropia de Shannon e PageRank Bipartido)}
No grafo bipartido $G_{\text{bip}}$, para cada unidade orçamentária $u \in U$ com orçamento empenhado total $W_u = \sum_{c \in N(u)} w_{u,c}$ e participação relativa do credor $s_{u,c} = w_{u,c} / W_u$, calcula-se o Índice Herfindahl--Hirschman ($\text{HHI}_u$) em $\mathcal{O}(|E|)$~\cite{rhoades1993herfindahl}:
\begin{equation}
  \text{HHI}_u = 10.000 \sum_{c \in N(u)} s_{u,c}^2 \in (0,\, 10.000].
\end{equation}
Seguindo as faixas referenciais do DOJ/FTC~\cite{doj2010horizontal}, valores abaixo de $1.500$ indicam baixa concentração, de $1.500$ a $2.499$ concentração moderada e $\ge 2.500$ alta concentração. Para detectar dispersão versus dependência em poucos fornecedores, calcula-se a Entropia de Shannon Normalizada~\cite{shannon1948mathematical} sobre a distribuição orçamentária da unidade ($N_u = |N(u)| > 1$):
\begin{equation}
  H_{\text{norm}}(u) = \frac{-\sum_{c \in N(u)} s_{u,c} \ln(s_{u,c})}{\ln(N_u)} \in [0,\, 1],
\end{equation}
em que valores próximos de zero denotam extrema concentração na unidade $u$. Aplicando contração de vértices sobre credores ligados ao mesmo componente societário do Eixo 1, computa-se $\Delta\text{HHI}_u = \text{HHI}_{u,\text{grupo}} - \text{HHI}_u$. Por fim, o algoritmo \textit{PageRank} ponderado~\cite{brin1998anatomy} ($\mathcal{O}(|V|+|E|)$ por iteração) ranqueia a relevância sistêmica conjunta das secretarias concentradoras e dos credores que atuam como pontes entre múltiplas pastas municipais.
